% !TeX root = ../main.tex

\section{Nutzeranforderungen}
\label{appsec:nutzer-anforderungen}

Dieser Abschnitt beschreibt knapp und für alle Stakeholder verständlich, welche
Anforderungen an das System gestellt werden.

\input{helper/requirements/content/user-requirements}

\section{Funktionale Anforderungen}
\label{appsec:funktionale-anforderungen}

Dieser Abschnitt beschreibt, welche Funktionen das System bereitstellen muss, wie es auf
bestimmte Eingaben reagiert und wie es sich in bestimmten Situationen verhalten soll.

\input{helper/requirements/content/functional-requirements}

\section{Nicht-funktionale Anforderungen}
\label{appsec:nicht-funktionale-anforderungen}

Dieser Abschnitt beschreibt, welche Rahmenbedingungen für das System gelten, z.\,B. in
Bezug auf Leistung, Sicherheit oder Benutzerfreundlichkeit.

\input{helper/requirements/content/non-functional-requirements}

\section{Anforderungsübersicht}
\label{appsec:anforderungsuebersicht}

Diese Tabelle dient als Übersicht über alle in den obigen Abschnitten definierten
Anforderungen.

\input{helper/requirements/content/requirements-overview}
